\documentclass[a4paper,12pt]{article}
\usepackage[utf8]{inputenc}
\usepackage[ngerman]{babel}
\usepackage{amsmath}
\usepackage{parskip}

\setlength{\parindent}{0pt}

\title{Modul 293 - Aufgabenblatt 12}
\author{von Lukas Meier}
\date{Unterricht vom 24.03.2026}

\begin{document}

\maketitle

\section*{Aufgabenblatt: HTML-Formulare II}

\subsection*{Teil 1: Elemente zuordnen}
Ordne die passenden Formularelemente zu:
\begin{itemize}
  \item Auswahl genau einer Option aus mehreren
  \item Auswahl mehrerer unabhängiger Optionen
  \item Auswahl aus einer Liste
  \item längerer Freitext
\end{itemize}

Verwende dabei die Begriffe \texttt{select}, \texttt{option}, \texttt{radio}, \texttt{checkbox}, \texttt{textarea}.

\subsection*{Teil 2: Dropdown bauen}
Erstelle ein Formular mit einem Dropdown zur Auswahl eines Kurses.

Anforderungen:
\begin{itemize}
  \item \texttt{select}-Element
  \item mindestens vier \texttt{option}-Einträge
  \item sinnvolles \texttt{name}-Attribut
\end{itemize}

\subsection*{Teil 3: Radios und Checkboxen}
Erweitere das Formular aus Teil 2 mit:
\begin{itemize}
  \item Radio-Buttons für den Schwierigkeitsgrad: Anfänger, Mittel, Fortgeschritten
  \item Checkboxen für Zusatzwünsche: Pausenkaffee, Handout, Zertifikat
\end{itemize}

\subsection*{Teil 4: Textarea ergänzen}
Füge eine \texttt{textarea} hinzu.
Darin sollen Teilnehmende Bemerkungen eingeben können.

Beantworte zusätzlich:
\begin{itemize}
  \item Worin unterscheidet sich \texttt{textarea} von einem normalen Textfeld?
\end{itemize}

\subsection*{Teil 5: Code lesen}
Betrachte den folgenden Code:

\begin{verbatim}
<form action="/feedback" method="post">
  <select name="thema">
    <option>HTML</option>
    <option>CSS</option>
  </select>

  <input type="radio" name="tempo" value="langsam">
  <input type="radio" name="tempo" value="schnell">

  <input type="checkbox" name="material" value="handout">
  <textarea name="kommentar"></textarea>
</form>
\end{verbatim}

\begin{itemize}
  \item Welche Arten von Eingaben sind hier möglich?
  \item Warum haben die Radio-Buttons denselben \texttt{name}-Wert?
  \item Welche Eingaben könnten mehrfach gewählt werden?
\end{itemize}

\subsection*{Teil 6: Praxisaufgabe}
Erstelle eine kleine HTML-Seite mit einem Feedback- oder Anmeldeformular.

Anforderungen:
\begin{itemize}
  \item ein Dropdown
  \item mindestens drei Radio-Buttons
  \item mindestens zwei Checkboxen
  \item eine \texttt{textarea}
  \item ein Submit-Button
\end{itemize}

\subsection*{Teil 7: Reflexion}
Beantworte kurz:
\begin{itemize}
  \item Wann verwendest du ein Dropdown statt Radio-Buttons?
  \item Wann sind Checkboxen sinnvoll?
  \item Wo hattest du heute noch Unsicherheiten?
\end{itemize}

\end{document}
